\section{Motivating Example}
\label{sec:motivation}

\subsection{Example: Generating Tests for P4 Implementations}
\label{sec:testing-type-checking}

Consider two properties of P4 implementations within a stated language fragment and target: well-typed programs should be accepted, and ill-typed programs should be rejected.
P4Smith~\cite{ruffy20gauntlet} generates valid P4 programs, while Failing with Purpose~\cite{lee26fwp} tests rejection by mutating well-typed programs into ill-typed ones guided by dangling coverage.

These examples raise three questions about the resulting tests.
First, are the properties' input conditions and expected outcomes precisely stated and justified over the specification?
Second, can each generator produce all and only the inputs satisfying its property's precondition---well-typed programs for the positive case and ill-typed programs for the negative case?
For either approach, producing suitable samples alone does not establish that the set of possible outputs matches the set defined by the precondition.
Third, how should the implementation be run and its observations interpreted to evaluate the property's goal, distinguishing acceptance and reported rejection from crashes or timeouts?

\subsection{Three Challenges for Trustworthy PBT}
\label{sec:three-propositions}

\paragraph{Formalizing and proving properties.}
\label{sec:property-proposition}
The intended property must have a precise mathematical statement and a proof over the specification that makes its assumptions and scope explicit.

\paragraph{Constructing and verifying generators.}
\label{sec:generator-propositions}
The generator must be sound and complete for the property's precondition: it produces only inputs satisfying that precondition and can produce every such input.
Soundness and completeness concern this set of possible outputs; generation efficiency and sampling diversity must be assessed separately.

\paragraph{Testing external implementations.}
\label{sec:implementation-procedure}
Testing implementations requires a procedure that, for any property, determines what to run on each generated input and which observations of the result to check against the property's goal.

Section~\ref{sec:instantiation} addresses these challenges using SpecTec and Lean: property formalization and proof in Section~\ref{sec:propose}, generator construction and verification in Section~\ref{sec:generators}, and implementation testing in Section~\ref{sec:test}.
